% Part: methods
% Chapter: proofs
% Section: cant-do-it

\documentclass[../../../include/open-logic-section]{subfiles}

\begin{document}

\olfileid{mth}{prf}{cnt}

\olsection{میرے کولوں نہیں ہُندا!{}}

اسیں سارے کدی اوس تھاں پہنچ جاندے آں جتھے کم چھڈن نوں جی کردا اے۔ پر تسی ایہ کر \emph{سکدے} او۔ تہاڈے استاد، پڑھاؤن والے معاون تے جماعت دے ساتھی مدد کر سکدے نیں۔ اوہناں توں مدد منگو!{} ہِیٹھ کجھ گلّاں نیں جیہڑیاں اوکڑ توں بچن وچ مدد کرن گئیاں تے دسن گئیاں کہ جے کم چھڈن نوں جی کرے تاں کی کریے۔

مسئلے حل کر سکّن لئی ایہ کم کرو:
\begin{enumerate}
\item \emph{جِنّا ہو سکے پہلاں کم شروع کرو۔} سمسٹر وچ اسیں رُجھے رہندے آں تے ساڈے وچوں کئی کم اگے پاؤندے رہندے آں۔ تسی جیہڑے سبھ توں چنگے کم کر سکدے او اوہناں وچوں اک ایہ اے کہ گھر دا کم سویرے شروع کرو۔ ایویں جے پھس جاؤ، تاں حل لبھن لئی تہاڈے کول ویلا ہووے گا (رون توں علاوہ)۔
\item \emph{اپنی جماعت دے ساتھیاں نال گل کرو}۔ تسی اکلّے نہیں۔ جماعت وچ ہوراں نوں وی اوکڑ ہو سکدی اے، پر شاید وکھ وکھ گلاں وچ۔ ساتھیاں نال گل کرن نال مسئلہ کسے ہور پاسے توں دِس سکدا اے تے حل دا راہ کھل سکدا اے۔ بے شک اوہناں دا حل بس نقل نہ کرو۔ اوہناں توں اشارہ منگو، یا دسو کہ کتھے پھسے او تے اگلا قدم پُچھو۔ تے جدوں تہانوں سمجھ آ جائے، تسی وی اوہناں دی مدد کرو۔ کسے ہور نوں اگے ودھاؤن تے گل سمجھاؤن نال تہانوں آپ وی ہور سمجھ آوے گی۔
\item \emph{مدد منگو۔} مدد لئی تہاڈے کول کئی وسیلے نیں۔ تہاڈے استاد تے پڑھاؤن والے معاون تہاڈی مدد لئی نیں تے \emph{چاہندے} نیں کہ تسی کامیاب ہووو۔ اوہ تہانوں مسئلہ حل کرن وچ مدد دے سکن تے لبھ سکن کہ تہانوں کیہڑے قدم اُتے اوکڑ آؤندی اے۔
\item \emph{تھوڑا آرام کرو۔} جے تسی پھسے او، تاں \emph{ہو سکدا اے} تسی بڑی دیر توں مسئلے ول تکّ رہے او۔ تھوڑی دیر آرام کرو، چاء دا اک کپ پیو، یا کجھ دیر کسے ہور مسئلے اُتے کم کرو، تے فیر تازہ ذہن نال ایس مسئلے ول مُڑو۔ اک رات سَو کے فیر ویکھ لو۔
\end{enumerate}

دھیان دتا کہ ایہناں طریقیاں لئی ثبوت اُتے کم کافی پہلاں شروع کرنا ضروری اے؟ جے تسی جمع کراؤن توں پہلی رات دو وجے ثبوت شروع کیتا، تاں شاید ایہ اینے کم نہ آون۔

ایہ گل شاید بڑی مایوسی والی لگے، پر ثبوت لبھنا اک اوکھا کم اے جیہدا پھل آخر ملدا اے۔ کجھ لوک ایہ کم پیشے وجوں کردے نیں، سو ایہدے وچ کوئی سواد ضرور ہووے گا۔ لگ بھگ ہر کم وانگ مسئلے حل کرن تے ثبوت دین لئی مشق چاہیدی اے۔ تہاڈی جماعت وچ کجھ ساتھیاں نوں ایہ کم سوکھا لگ سکدا اے۔ اوہناں نے شاید پہلاں ای بہت مشق کیتی اے۔ چھیتی ہار نہ منّو۔

جے کسے مسئلے لئی تہاڈا ویلا (یا صبر) مُک جائے، تاں کوئی گل نہیں۔ ایس دا مطلب ایہ نہیں کہ تسی بے وقوف او یا تہانوں کدی سمجھ نہیں آوے گی۔ (اپنے استاد یا کسے ہور طالب علم توں) پتا کرو کہ کم کیویں ہُندا اے، تے لبھو کہ کتھے غلطی ہوئی یا کتھے پھسے سو، تاں جو اگلی واری ایہو جیہی اوکڑ آوے تاں اوہ گل فیر نہ ہووے۔ فیر حل ویکھے بناں ایہ کم کرن دی کوشش کرو۔ تے اگلی واری پہلاں شروع کرو تے پہلاں ای مدد منگو۔

\end{document}
